\documentclass[11pt,a4paper]{article}
\usepackage[T1]{fontenc}
\usepackage{lmodern,amsmath,amssymb,booktabs,microtype}
\usepackage[margin=25mm]{geometry}
\usepackage[colorlinks=true,urlcolor=blue,linkcolor=blue]{hyperref}
\newcommand{\tr}{\operatorname{tr}}
\newcommand{\Tr}{\operatorname{Tr}}
\newcommand{\E}{\mathbb E}
\newcommand{\HH}{\mathcal H}
\title{A dynamical singlet moment in complex DSSYK\\Gaussian contractions and a symmetry reduction}
\author{Alan Day \and Chengheng\thanks{AI research assistant (OpenAI).}}
\date{Local research note --- 5 November 2026}
\begin{document}\maketitle
\begin{abstract}
The static-patch programme needs singlet dynamics that can be compared with
bulk propagation. We organize the first derivative singlet's equal-time
strength and second frequency moment as exact Gaussian contractions.
Applying the established multiplicity-free decomposition of fermionic
operator spaces eliminates the linear Gaussian component of its time
derivative. The dynamical moment consequently reduces to a positive cubic
tensor norm. A complete finite-system contraction supplies a reproducible
benchmark for the subsequent double-scaled calculation.
\end{abstract}

\section{The physical quantity to calculate}
Our target remains the singlet-to-bulk-particle step of the Susskind
programme \cite{mss}. In the conventions of the preceding dictionary note,
\begin{equation}
 H=\sum_{I,J}G_{IJ}c_I^\dagger c_J,\quad
 \E G_{IJ}G_{KL}=w\delta_{IL}\delta_{JK},\quad
 w=4^p\frac{\mathcal J_s^2}{\lambda}
             \left(\frac{p!}{N^p}\right)^2,\quad \lambda=\frac{p^2}{N}.
\end{equation}
Here $G$ is Hermitian, each index set has size $p$, the CAR is standard,
and $H=H_s$ is in string units. At infinite temperature
$\tr=2^{-N}\Tr$. Define
\begin{equation}
 A=\sum_i[H,c_i]^\dagger[H,c_i],\qquad K=[H,A].
\end{equation}
With $C(t)=\E\tr(A(t)A)$ and $S(\omega)=\int dt\,e^{i\omega t}C(t)$,
\begin{equation}
 \mu_0=\E\tr A^2,\qquad
 \boxed{\mu_2=\int\frac{d\omega}{2\pi}\omega^2S(\omega)
             =\E\tr K^\dagger K=-C''(0).}\label{moment}
\end{equation}
The $\omega^2$ weight removes every zero-frequency contribution. Thus
$\mu_2$ is a direct measure of singlet dynamics even when the equal-time
strength is dominated by conserved components. The next physical target is
its double-scaled size, followed by enough further moments to resolve the
frequency dependence. A mass assignment will then require the bulk
propagation kernel, rather than identifying a frequency moment with a mass.

\newpage
\section{A contraction formula that keeps the same couplings}
Choose real standard Gaussian coordinates $x_a$ for the GUE and absorb
their variances into Hermitian matrices $V_a$, so $H=\sum_a x_aV_a$ and
$\E x_ax_b=\delta_{ab}$. The same coordinates occur in $A$:
\begin{align}
 D_{ai}&=[V_a,c_i],\\
 B_{ab}&=\frac12\sum_i(D_{ai}^\dagger D_{bi}+D_{bi}^\dagger D_{ai}),
 & A&=\sum_{ab}x_ax_bB_{ab},\\
 C_{abc}&=\frac13\bigl([V_a,B_{bc}]+[V_b,B_{ac}]+[V_c,B_{ab}]\bigr),
 & K&=\sum_{abc}x_ax_bx_c C_{abc}.
\end{align}
$B$ is symmetric and Hermitian; $C$ is fully symmetric and
anti-Hermitian. Write $\|X\|^2=\tr X^\dagger X$, and set
\begin{equation}
 \overline A=\sum_a B_{aa},\qquad L_c=\sum_a C_{aac}.
\end{equation}
Ordinary Gaussian Wick contraction, with no large-$N$ approximation, gives
\begin{align}
 \mu_0&=\|\overline A\|^2+2\sum_{ab}\|B_{ab}\|^2,\label{mu0}\\
 \mu_2&=6\sum_{abc}\|C_{abc}\|^2+9\sum_c\|L_c\|^2.\label{mu2general}
\end{align}
The factors $6$ and $9$ count, respectively, the pairings with three
cross-contractions and the pairings with one internal contraction in each
cubic factor. Equivalently,
\begin{equation}
 K=\sum_{abc}C_{abc}:x_ax_bx_c:+3\sum_cL_cx_c,
\end{equation}
where the colons denote ordinary Gaussian normal ordering of the
\emph{couplings}. Orthogonality of these two Gaussian degrees proves
\eqref{mu2general} without subtracting large signed contributions.

\paragraph{Symmetry reduction.}
For the full complex-SYK GUE ensemble,
\begin{equation}
 \boxed{L_c=0,\qquad \mu_2=6\sum_{abc}\|C_{abc}\|^2.}\label{reduction}
\end{equation}
Appendix A proves this by applying the known multiplicity-free $U(N)$
operator decomposition \cite{kota}. The reduction concerns degree in
Gaussian couplings, not a statement about a number of propagating bulk
particles. It holds sector by sector before taking a large-$N$ limit.

\newpage
\section{A complete finite-system benchmark}
For $N=4,p=2$, a full covariance sum gives the following rational values.
The table uses $w=1$; multiply the last two columns by $w^2,w^3$ to restore
the physical variance. Each sector trace is divided by $d_m=\binom Nm$.
\begin{center}
\begin{tabular}{rrrrr}\toprule
$m$ & $d_m$ & $\E\tr_m A$ & $\E\tr_m A^2$ & $\E\tr_m K^\dagger K$\\\midrule
0 & 1 & 0 & 0 & 0\\
1 & 4 & 0 & 0 & 0\\
2 & 6 & 12 & 292 & 0\\
3 & 4 & 27 & 942 & 480\\
4 & 1 & 36 & 1488 & 0\\\midrule
full & 16 & $27/2$ & 438 & 120\\\bottomrule
\end{tabular}
\end{center}
The only dynamical sector here is $m=3$. At $\mathcal J_s=1$ one has
$w=1/4$, hence
\begin{equation}
 \mu_0=\frac{219}{8},\qquad \boxed{\mu_2=\frac{15}{8}}.
\end{equation}
This positive value is an ensemble moment, not a random-sample estimate.
The $m=p$ block is a useful physical control: $A_p=pH_p^2$ is conserved.
The full $p=1$ model is another control, with
$A=c^\dagger h^2c$ for $H=c^\dagger hc$ and $\mu_2=0$.

An independent expansion of the unsymmetrized cubic expression into all
15 Wick pairings yields nine zero terms, two terms $645$, and four terms
$-585/2$, summing to $120$. Their absolute sum is $20.5$ times the
answer; the positive expression \eqref{reduction} avoids this cancellation.
The integer contraction $L_c$ vanishes exactly in the finite calculation.

As an additional check, a three-coordinate restriction of this same
matrix model breaks the full ensemble symmetry. Both Gaussian degrees
then contribute: $\mu_2^{(3)}=1/8$ and $\mu_2^{(1)}=1/32$.
A 64-node product Gaussian quadrature, exact through degree six, agrees
with their sum to scaled error below $10^{-15}$. This checks the term
eliminated by symmetry as well as the surviving term.

\section{The next analytic step}
Existing $q$-Hermite and chord techniques solve substantial parts of
embedded-ensemble spectral theory \cite{tt}. Here we first reduce the
finite-$N$ tensor norm to body-rank sums using the covariance eigenchannels
of \cite{kota}. The target is the leading dependence on $N,p$ at fixed
$p^2/N$, then its small positive-$\lambda$ expansion. The benchmark above
fixes the normalization and signed contractions of that calculation.

\newpage
\appendix
\section{Why the linear Gaussian component vanishes}
This is an application of the standard fermionic operator decomposition,
not a new representation-theoretic theorem. Let
$\HH_r=\bigwedge^r\mathbb C^N$ and let $\mathcal E_m$ embed a $p$-particle
coupling matrix into the $m$-particle Hamiltonian. The known decomposition
\cite{kota}, Section II, is
\begin{equation}
 \operatorname{End}(\HH_r)=
 \bigoplus_{s=0}^{\min(r,N-r)}\mathcal R_s.
\end{equation}
Each $\mathcal R_s$ occurs once; different $s$ are inequivalent. The
summands are preserved by adjoint and entrywise complex conjugation.
Normal-ordered embedding $\mathcal E_m$ commutes with both operations
and is nonzero on every irrep common to source and target for $m\ge p$.
For example its highest-weight rank-$s$ monomial is multiplied by a
nonzero real spectator factor proportional to $\binom{m-s}{p-s}$.
For $m<p$, both $H_m$ and $A_m$ vanish.

For a Hermitian source matrix $X$, form the linear map
\begin{equation}
 \mathcal F_m(X)=\E\left[K_m(G)\,\Tr_{\HH_p}(GX)\right].
\end{equation}
The GUE measure and the sum over fermion labels in $A$ are $U(N)$
covariant. Thus $\mathcal F_m$ is an equivariant real-linear map on
Hermitian matrices, with a complex-linear extension to their full matrix
space. Multiplicity freedom implies that on each common irrep
\begin{equation}
 \mathcal F_m|_{\mathcal R_s}=\alpha_s\mathcal E_m|_{\mathcal R_s}.
\end{equation}
All CAR matrices may be chosen real in the occupation basis. Conjugating
$G$ therefore conjugates $K$, and the GUE law is unchanged. Consequently
$\mathcal F_m(\overline X)=\overline{\mathcal F_m(X)}$, which makes
$\alpha_s$ real. On the other hand $K_m$ is anti-Hermitian, whereas
$\mathcal E_m(X)$ is Hermitian for Hermitian $X$; this makes $\alpha_s$
purely imaginary. Hence $\alpha_s=0$. Irreps absent from the target have
zero image by equivariance. We conclude $\mathcal F_m=0$.

Every real Gaussian coordinate is a real linear functional of $G$.
It follows that $\E[x_cK_m]=0$. Ordinary Gaussian integration gives
$\E[x_cK_m]=3\sum_a(C_{aac})_m$, proving \eqref{reduction} in every
sector. A restricted set of GUE coordinates is not invariant under
$U(N)$, which explains why the independent three-coordinate test retains
a nonzero linear component.

\section{Conventions, verification and scope}
The calculation uses finite $N$, $1\le p\le N$, the full Hermitian
Gaussian covariance, $\beta=\mu=0$, and no spectral truncation. It is
ordinary Wick ordering in coupling space; it is not an identification
with the $q$-Hermite normal ordering of the limiting Hamiltonian.
For $H_c=pH_s$, $A_c=p^2A_s$, the corresponding moments obey
$\mu_{0,c}=p^4\mu_{0,s}$ and $\mu_{2,c}=p^6\mu_{2,s}$.

\newpage
The tensor sum has not yet been reduced to a closed combinatorial
expression at general $N,p$. Its present matrix implementation scales
poorly and is a calibration tool, not the proposed large-$N$ algorithm.
A nonzero finite-$N$ moment alone does not establish its double-scaled
limit or a particle mass. Subtracting $\E A$ or the sample trace of $A$
also does not remove every conserved operator. Equation \eqref{moment}
is insensitive to all such zero-frequency contributions.

\paragraph{Reproducibility.}
\texttt{code/singlet\_wick\_moments.py} constructs all coupling directions,
the quadratic and cubic tensors, and their sector norms.
\texttt{code/run\_singlet\_wick\_moments.py} records full contractions
for $(N,p)=(3,1),(3,2),(4,2)$, the 15-pairing check, independent
quadrature, and a spectral-moment comparison. Integer-valued matrix
operations are checked before rational normalization; the control
quadrature and diagonalization use floating point. Parameters, dependency
versions, exact fractions and code hashes are stored in
\texttt{moments.json} in the corresponding results directory.
No disorder samples are added to the previous stage.

\paragraph{Literature use.}
We reuse the covariance decomposition of \cite{kota} and the Wick-method
perspective of \cite{tt}. The latter's independent-probe final kernel
does not itself evaluate the coupling-dependent $A$. The particle-removal
moments in \cite{kv} are a further route to the adjacent-sector sums and
are the next reading target. The original coupled four-point density
from the dictionary note remains the all-frequency object to solve.

\begin{thebibliography}{9}
\bibitem{mss} S. Miyashita, Y. Sekino and L. Susskind,
\href{https://arxiv.org/html/2506.18054v3}{arXiv:2506.18054v3}.
\bibitem{kota} V. K. B. Kota, \emph{SU(N) Wigner--Racah algebra for the
matrix of second moments of embedded Gaussian unitary ensemble of
random matrices}, J. Math. Phys. \textbf{46}, 033514 (2005),
\href{https://arxiv.org/html/nucl-th/0410069v1}{arXiv:nucl-th/0410069v1}.
\bibitem{tt} J. Tall and S. Tomsovic,
\href{https://arxiv.org/html/2604.14522v2}{arXiv:2604.14522v2},
Sections 2--4 and Appendix A.
\bibitem{kv} V. K. B. Kota and M. Vyas, \emph{Random Matrix Theory for
Transition Strength Densities in Finite Quantum Systems: Results from
Embedded Unitary Ensembles},
\href{https://arxiv.org/html/1504.00759v1}{arXiv:1504.00759v1}.
\end{thebibliography}
\end{document}
